\textbf{Phần 1. Trắc nghiệm lựa chọn}
\setcounterpageref{ex}{0}

\begin{ex} %[3D2V09]
Tính giá trị của biểu thức: $A = 6 	imes 7 - 48 : 6 + 15$.
\choice
{$45$}
{$51$}
{\True $49$}
{$47$}
\loigiai{
Thực hiện nhân chia trước, cộng trừ sau:\
$$A = 42 - 8 + 15 = 34 + 15 = 49$$
\textbf{Đáp án của câu hỏi là C}
}
\end{ex}

\begin{ex} %[3D2V09]
Một con kiến có $6$ cái chân. Một đàn kiến gồm $8$ con kiến thợ và $4$ con kiến chúa. Hỏi đàn kiến đó có tất cả bao nhiêu cái chân?
\choice
{\True $72$ cái chân}
{$68$ cái chân}
{$76$ cái chân}
{$64$ cái chân}
\loigiai{
Tổng số con kiến trong đàn là: $8 + 4 = 12$ (con).\
Đàn kiến có tất cả số cái chân là: $12 \times 6 = 72$ (cái chân).
\textbf{Đáp án của câu hỏi là A}
}
\end{ex}

\begin{ex} %[3D2V09]
Tìm giá trị của $x$, biết: $x : 6 + 25 = 32$.
\choice
{$36$}
{$48$}
{$30$}
{\True $42$}
\loigiai{
Ta có: $x : 6 = 32 - 25 = 7 \Rightarrow x = 7 \times 6 = 42$.
\textbf{Đáp án của câu hỏi là D}
}
\end{ex}

\begin{ex} %[3D2V09]
Một hình lục giác đều có 6 cạnh bằng nhau. Biết chu vi của hình lục giác đều đó là $54\text{ cm}$. Độ dài mỗi cạnh của hình lục giác đều là
\choice
{$8\text{ cm}$}
{\True $9\text{ cm}$}
{$7\text{ cm}$}
{$6\text{ cm}$}
\loigiai{
Vì hình lục giác đều có 6 cạnh bằng nhau nên độ dài mỗi cạnh là: $54 : 6 = 9\text{ (cm)}$.
\textbf{Đáp án của câu hỏi là B}
}
\end{ex}

\begin{ex} %[3D2V09]
Một cửa hàng có $60\text{ kg}$ gạo, người ta đóng vào các túi loại $6\text{ kg}$. Sau đó cửa hàng bán đi $4$ túi gạo. Hỏi cửa hàng còn lại bao nhiêu ki-lô-gam gạo?
\choice
{$24\text{ kg}$}
{$40\text{ kg}$}
{\True $36\text{ kg}$}
{$30\text{ kg}$}
\loigiai{
Cửa hàng đóng được số túi gạo là: $60 : 6 = 10$ (túi).\
Sau khi bán 4 túi, cửa hàng còn lại: $10 - 4 = 6$ (túi).\
Số ki-lô-gam gạo còn lại là: $6 \times 6 = 36\text{ (kg)}$.
\textbf{Đáp án của câu hỏi là C}
}
\end{ex}

\begin{ex} %[3D2V09]
Tính nhanh giá trị của biểu thức: $6 \times 2 + 6 \times 3 + 6 \times 5$.
\choice
{\True $60$}
{$54$}
{$66$}
{$72$}
\loigiai{
Đặt thừa số chung: $6 \times (2 + 3 + 5) = 6 \times 10 = 60$.
\textbf{Đáp án của câu hỏi là A}
}
\end{ex}

\begin{ex} %[3D2V09]
Trong một phép chia cho $6$, thương bằng $8$ và số dư là số dư lớn nhất có thể có. Số bị chia là
\choice
{$48$}
{$51$}
{$52$}
{\True $53$}
\loigiai{
Vì số chia là 6 nên số dư lớn nhất là 5.\
Số bị chia là: $6 \times 8 + 5 = 48 + 5 = 53$.
\textbf{Đáp án của câu hỏi là D}
}
\end{ex}

\begin{ex} %[3D2V09]
Điền dấu thích hợp vào chỗ chấm: $6 \times 8 - 14 \dots 54 : 6 \times 4$.
\choice
{$>$}
{\True $<$}
{$=$}
{Không so sánh được}
\loigiai{
Vế trái: $6 \times 8 - 14 = 48 - 14 = 34$.\
Vế phải: $54 : 6 \times 4 = 9 \times 4 = 36$.\
Vì $34 < 36$ nên điền dấu $<$. 
\textbf{Đáp án của câu hỏi là B}
}
\end{ex}

\begin{ex} %[3D2V09]
Một tổ công nhân sản xuất được $48$ chi tiết máy trong một tuần và chia đều thành 6 ca làm việc. Mỗi ca lại chia đều sản phẩm cho 4 công nhân. Hỏi mỗi công nhân làm được bao nhiêu chi tiết máy?
\choice
{$3$ chi tiết máy}
{$4$ chi tiết máy}
{\True $2$ chi tiết máy}
{$1$ chi tiết máy}
\loigiai{
Mỗi ca làm được: $48 : 6 = 8$ chi tiết máy.\
Mỗi công nhân làm được: $8 : 4 = 2$ chi tiết máy.
\textbf{Đáp án của câu hỏi là C}
}
\end{ex}

\begin{ex} %[3D2V09]
Tìm $x$ biết: $(x - 12) \times 6 = 42$.
\choice
{\True $19$}
{$18$}
{$20$}
{$21$}
\loigiai{
Ta có: $x - 12 = 42 : 6 = 7 \Rightarrow x = 7 + 12 = 19$.
\textbf{Đáp án của câu hỏi là A}
}
\end{ex}

\begin{ex} %[3D2V09]
Cho dãy số: $6; 12; 18; 24; \dots$ Số hạng thứ 9 của dãy số trên là
\choice
{$48$}
{$60$}
{$56$}
{\True $54$}
\loigiai{
Dãy số gồm các bội số của 6. Số hạng thứ 9 là: $6 \times 9 = 54$.
\textbf{Đáp án của câu hỏi là D}
}
\end{ex}

\begin{ex} %[3D2V09]
Mai có $18$ quyển vở, số vở của Lan gấp $6$ lần số vở của Mai. Hỏi cả hai bạn có tất cả bao nhiêu quyển vở?
\choice
{$108$ quyển vở}
{\True $126$ quyển vở}
{$120$ quyển vở}
{$114$ quyển vở}
\loigiai{
Số vở của Lan là: $18 \times 6 = 108$ (quyển).\
Cả hai bạn có: $18 + 108 = 126$ (quyển).
\textbf{Đáp án của câu hỏi là B}
}
\end{ex}

\textbf{Phần 2. Trắc nghiệm đúng sai}
\setcounterpageref{ex}{0}

\begin{ex} %[3D2V09]
Xét tính đúng sai của các khẳng định về phép nhân và phép chia với 6:
\choiceTF
{\True Tích của số 6 với một số tự nhiên bất kì luôn luôn là một số chẵn.}
{Trong phép chia cho 6, số dư lớn nhất có thể là 6.}
{\True Một số tự nhiên chia hết cho 6 thì chắc chắn chia hết cho 3.}
{\True Giá trị của $6 \times 4$ bằng giá trị của $4 \times 6$.}
\loigiai{
\begin{itemize}
    \item Ý a) Đúng: Vì 6 là số chẵn nên $6 \times n$ luôn là số chẵn.
    \item Ý b) Sai: Số dư luôn bé hơn số chia, lớn nhất là 5.
    \item Ý c) Đúng: Vì $6 = 2 \times 3$ nên số chia hết cho 6 cũng chia hết cho 3.
    \item Ý d) Đúng: Theo tính chất giao hoán của phép nhân.
\end{itemize}
\begin{itemize}
    \item a) \textbf{Đúng}.
    \item b) \textbf{Sai}.
    \item c) \textbf{Đúng}.
    \item d) \textbf{Đúng}.
\end{itemize}
}
\end{ex}

\begin{ex} %[3D2V09]
Một tấm vải dài $54\text{ m}$ được cắt thành các mảnh nhỏ dài $6\text{ m}$:
\choiceTF
{\True Người ta cắt được tất cả 9 mảnh vải nhỏ.}
{Người thợ phải thực hiện đúng 9 nhát cắt.}
{\True Tổng độ dài của 5 mảnh vải nhỏ là $30\text{ m}$.}
{\True Nếu bớt ở tấm vải đi $6\text{ m}$ thì số mảnh vải cắt được là 8 mảnh.}
\loigiai{
\begin{itemize}
    \item Ý a) Đúng: $54 : 6 = 9$ mảnh.
    \item Ý b) Sai: Để cắt thành 9 mảnh chỉ cần $9 - 1 = 8$ nhát cắt.
    \item Ý c) Đúng: $5 \times 6 = 30\text{ m}$.
    \item Ý d) Đúng: $(54 - 6) : 6 = 48 : 6 = 8$ mảnh.
\end{itemize}
\begin{itemize}
    \item a) \textbf{Đúng}.
    \item b) \textbf{Sai}.
    \item c) \textbf{Đúng}.
    \item d) \textbf{Đúng}.
\end{itemize}
}
\end{ex}

\textbf{Phần 3. Trả lời ngắn}
\setcounterpageref{ex}{0}

\begin{ex} %[3D2V09]
Tính giá trị của biểu thức: $6 \times 8 - 30 : 6$.
\shortans{43}
\loigiai{
\textbf{Đáp số:} $43$\\
$6 \times 8 - 30 : 6 = 48 - 5 = 43$.
}
\end{ex}

\begin{ex} %[3D2V09]
Một đàn ong có $7$ con. Hỏi đàn ong đó có tất cả bao nhiêu cái chân?
\shortans{42}
\loigiai{
\textbf{Đáp số:} $42$\\
Mỗi con ong có 6 chân: $7 \times 6 = 42$ cái chân.
}
\end{ex}

\begin{ex} %[3D2V09]
Tìm số bị chia trong phép chia cho 6, biết thương bằng 9 và số dư bằng 4.
\shortans{58}
\loigiai{
\textbf{Đáp số:} $58$\\
$6 \times 9 + 4 = 54 + 4 = 58$.
}
\end{ex}

\begin{ex} %[3D2V09]
Lớp 3A có 36 học sinh được chia đều vào 6 nhóm. Sau đó mỗi nhóm cử 2 bạn tham gia đá bóng. Hỏi có tất cả bao nhiêu bạn tham gia đá bóng?
\shortans{12}
\loigiai{
\textbf{Đáp số:} $12$\\
Mỗi nhóm có $36 : 6 = 6$ bạn. Có $6 \times 2 = 12$ bạn tham gia.
}
\end{ex}

\begin{ex} %[3D2V09]
Tìm giá trị của $x$, biết: $x : 6 = 12 - 4$.
\shortans{48}
\loigiai{
\textbf{Đáp số:} $48$\\
$x : 6 = 8 \Rightarrow x = 8 \times 6 = 48$.
}
\end{ex}

\begin{ex} %[3D2V09]
Tính nhanh: $6 \times 1 + 6 \times 2 + 6 \times 3 + 6 \times 4$.
\shortans{60}
\loigiai{
\textbf{Đáp số:} $60$\\
$6 \times (1 + 2 + 3 + 4) = 6 \times 10 = 60$.
}
\end{ex}
